On prépare une solution aqueuse $\mathrm{(S)}$ d'acide azélaïque (noté
$\ce{AH}$, et sa base conjuguée $\ce{A-}$) de concentration
$C=\qty{1.00e-2}{\M}$. La mesure de $\mathrm{pH}$ de la solution à
$\qty{25}{\degreeCelsius}$ donne $\mathrm{pH}=3,28$.

\begin{enumerate}
    \item Écrire l'équation qui modélise la réaction entre l'acide $\ce{AH}$ et
          l'eau.~\textbf{(1pt)}
    \item Dresser le tableau d'avancement de la réaction entre l'acide azélaïque
          et l'eau.~\textbf{(0,75pt)}
    \item Montrer que le taux d'avancement final de la réaction s'écrit
          $\tau=\frac{10^{-\mathrm{pH}}}{C}$.~\textbf{(1pt)}
    \item Calculer la valeur de $\tau$. Déduire.~\textbf{(0,5pt)}
    \item Montrer que l'expression du quotient de la réaction $Q_{r,\eq}$ à
          l'équilibre s'écrit
          $Q_{r,\eq}=\frac{C\cdot\tau^{2}}{1-\tau}$.~\textbf{(1pt)}
    \item Vérifier que la valeur du $\mathrm{pK_{A}}$ du couple
          $(\ce{AH_{(aq)}}\ttslash{}\ce{A^-_{(aq)}})$ est
          $\mathrm{pK_{A}}=4,54$.~\textbf{(1pt)}
    \item Déterminer, parmi les espèces $\ce{AH_{(aq)}}$ et $\ce{A^-_{(aq)}}$,
          l'espèce qui prédomine dans la solution à
          l'équilibre.~\textbf{(0,5pt)}
\end{enumerate}

